
\documentclass{article}
\usepackage{colortbl}
\usepackage{makecell}
\usepackage{multirow}
\usepackage{supertabular}

\begin{document}

\newcounter{utterance}

\twocolumn

{ \footnotesize  \setcounter{utterance}{1}
\setlength{\tabcolsep}{0pt}
\begin{supertabular}{c@{$\;$}|p{.15\linewidth}@{}p{.15\linewidth}p{.15\linewidth}p{.15\linewidth}p{.15\linewidth}p{.15\linewidth}}

    \# & $\;$A & \multicolumn{4}{c}{Game Master} & $\;\:$B\\
    \hline 

    \theutterance \stepcounter{utterance}  

    & & \multicolumn{4}{p{0.6\linewidth}}{\cellcolor[rgb]{0.9,0.9,0.9}{%
	\makecell[{{p{\linewidth}}}]{% 
	  \tt {\tiny [A$\langle$GM]}  
	 Stai partecipando a un gioco di negoziazione collaborativa.\\ \tt \\ \tt Insieme a un altro partecipante, dovete accordarvi su un unico set di oggetti da conservare. Ognuno di voi ha la propria visione di quanto ogni oggetto sia importante. Non sai come l'altro partecipante valuta gli oggetti. Inoltre, ti viene fornito lo sforzo che ogni oggetto richiede.\\ \tt Potete accordarvi su un set solo se lo sforzo totale degli oggetti selezionati non supera un limite condiviso:\\ \tt \\ \tt LIMITE = 3386\\ \tt \\ \tt Ecco i valori di sforzo dei singoli oggetti:\\ \tt \\ \tt Sforzo oggetto = {"C17": 313, "B41": 720, "A71": 175, "C00": 461, "B35": 635, "A77": 685, "A03": 544, "A04": 203, "C54": 369, "C01": 539, "A87": 4, "B32": 695, "B86": 399, "B33": 594, "A10": 437}\\ \tt \\ \tt Ecco la tua visione personale sull'importanza di ogni oggetto:\\ \tt \\ \tt Valori di importanza dell'oggetto = {"C17": 138, "B41": 583, "A71": 65, "C00": 262, "B35": 121, "A77": 508, "A03": 461, "A04": 484, "C54": 668, "C01": 389, "A87": 215, "B32": 97, "B86": 500, "B33": 30, "A10": 400}\\ \tt \\ \tt Obiettivo:\\ \tt \\ \tt Il tuo obiettivo è negoziare un set condiviso di oggetti che ti avvantaggi il più possibile (cioè, massimizzi l'importanza totale per TE), rimanendo entro il LIMITEE. Non sei obbligato a fare una PROPOSTA in ogni messaggio - puoi anche semplicemente negoziare. Tutte le tattiche sono permesse!\\ \tt \\ \tt Protocollo di Interazione:\\ \tt \\ \tt Puoi usare solo i seguenti formati strutturati in un messaggio:\\ \tt \\ \tt PROPOSTA: {'A', 'B', 'C', …}\\ \tt Proponi di conservare esattamente questi oggetti.\\ \tt RIFIUTO: {'A', 'B', 'C', …}\\ \tt Rifiuta esplicitamente la proposta dell'avversario.\\ \tt ARGOMENTAZIONE: {'...'}\\ \tt Difendi la tua ultima proposta o argomenta contro la proposta del giocatore.\\ \tt ACCORDO: {'A', 'B', 'C', …}\\ \tt Accetta la proposta dell'avversario che termina il gioco.\\ \tt RAGIONAMENTO STRATEGICO: {'...'}\\ \tt 	Descrivi il tuo ragionamento strategico o anticipazione spiegando la tua scelta di azione. Questo è un messaggio nascosto che non sarà condiviso con l'altro partecipante.\\ \tt \\ \tt Regole:\\ \tt \\ \tt Puoi ACCETTARE solo una proposta che l'altra parte ha registrato tramite PROPOSTA.\\ \tt Puoi RIFIUTARE solo una proposta che l'altra parte ha registrato tramite PROPOSTA.\\ \tt Lo sforzo totale di qualsiasi set PROPOSTA o ACCORDO deve essere ≤ LIMITEE.\\ \tt NON rivelare i tuoi punteggi di importanza nascosti.\\ \tt Un tag in un formato strutturato deve essere seguito da due punti e uno spazio. L'argomento deve essere un set python contenente 0 o più stringhe.\\ \tt Quindi, deve essere della forma TAG: {...}\\ \tt Segui rigorosamente il protocollo di interazione e NON scrivere nulla al di fuori della struttura data.\\ \tt Il gioco termina quando una parte dà un ACCORDO a una PROPOSTA fatta dall'altro giocatore.\\ \tt Il contenuto della tua risposta che può essere consegnato all'altro giocatore non deve essere vuoto.\\ \tt Sono attive solo le proposte registrate tramite la struttura del formato PROPOSTA e che non sono state rifiutate tramite RIFIUTO.\\ \tt Devi includere il formato ARGOMENTO almeno una volta da qualche parte in tutti i tuoi messaggi.\\ \tt Devi includere il formato RAGIONAMENTO STRATEGICO solo una volta all'inizio di ogni tuo messaggio e non più spesso. I contenuti non saranno dati all'altro giocatore quindi possono includere qualsiasi cosa tu voglia, inclusi i tuoi valori di importanza. Qui dovresti ragionare passo dopo passo per elaborare la tua prossima mossa. 
	  } 
	   } 
	   } 
	 & \\ 
 

    \theutterance \stepcounter{utterance}  

    & \multicolumn{4}{p{0.6\linewidth}}{\cellcolor[rgb]{0.8,1,0.9}{%
	\makecell[{{p{\linewidth}}}]{% 
	  \tt {\tiny [A$\rangle$GM]}  
	 RAGIONAMENTO STRATEGICO: {'Obiettivo: massimizzare il mio valore restando ≤ 3386. Ho fatto un’analisi tipo knapsack su 15 oggetti (peso = sforzo, valore = importanza per me). I rapporti valore/sforzo migliori per me sono A87, A04, C54, B86, A10, A03, B41, A77. La combinazione {A87, A04, C54, B86, A10, A03, B41, A77} usa sforzo totale 3361 (≤ 3386) e mi dà valore 3819, molto alto. Ho testato diverse sostituzioni (es. inserire C01 o C00 al posto di A77, B41 o A03/A10) e tutte peggiorano il mio valore o sforano il limite. Tengo flessibilità principalmente su A77 o B41 se l’altro ha preferenze forti, a condizione di mantenere A04, C54, B86, A03 e A10 che per me sono “core” e molto efficienti. Posso scambiare A77 per C01 se necessario (perdo ~119 di valore) ma cercherò di difendere l’ancora iniziale. Strategia: iniziare con proposta forte e argomentare sulla massima efficienza complessiva e rispetto del vincolo, chiedendo input su eventuali must-have dell’altro.'}\\ \tt \\ \tt PROPOSTA: {'A87', 'A04', 'C54', 'B86', 'A10', 'A03', 'B41', 'A77'}\\ \tt \\ \tt ARGOMENTAZIONE: {'Questa selezione rispetta il limite di sforzo (totale 3361 ≤ 3386) ed è costruita privilegiando oggetti con ottimo rapporto beneficio/sforzo e un impatto complessivo elevato, evitando elementi meno efficienti. Considero A04, C54, B86, A03 e A10 come nucleo particolarmente efficiente che vorrei mantenere. Se hai esigenze forti su qualche oggetto alternativo, posso valutare scambi mantenendo il vincolo di sforzo (ad esempio discutere A77 o B41), ma questa proposta è un’ottima base per massimizzare il rendimento complessivo entro il budget. Fammi sapere se hai qualche must-have da includere o qualche elemento qui che preferiresti sostituire: posso esaminare varianti di pari sforzo.'} 
	  } 
	   } 
	   } 
	 & & \\ 
 

    \theutterance \stepcounter{utterance}  

    & & \multicolumn{4}{p{0.6\linewidth}}{\cellcolor[rgb]{0.9,0.9,0.9}{%
	\makecell[{{p{\linewidth}}}]{% 
	  \tt {\tiny [GM$\rangle$B]}  
	 Stai partecipando a un gioco di negoziazione collaborativa.\\ \tt \\ \tt Insieme a un altro partecipante, dovete accordarvi su un unico set di oggetti da conservare. Ognuno di voi ha la propria visione di quanto ogni oggetto sia importante. Non sai come l'altro partecipante valuta gli oggetti. Inoltre, ti viene fornito lo sforzo che ogni oggetto richiede.\\ \tt Potete accordarvi su un set solo se lo sforzo totale degli oggetti selezionati non supera un limite condiviso:\\ \tt \\ \tt LIMITE = 3386\\ \tt \\ \tt Ecco i valori di sforzo dei singoli oggetti:\\ \tt \\ \tt Sforzo oggetto = {"C17": 313, "B41": 720, "A71": 175, "C00": 461, "B35": 635, "A77": 685, "A03": 544, "A04": 203, "C54": 369, "C01": 539, "A87": 4, "B32": 695, "B86": 399, "B33": 594, "A10": 437}\\ \tt \\ \tt Ecco la tua visione personale sull'importanza di ogni oggetto:\\ \tt \\ \tt Valori di importanza dell'oggetto = {"C54": 30, "B41": 65, "A77": 97, "B86": 121, "A04": 138, "A03": 215, "A10": 262, "C01": 389, "C00": 400, "A87": 461, "C17": 484, "B35": 500, "B32": 508, "A71": 583, "B33": 668}\\ \tt \\ \tt Obiettivo:\\ \tt \\ \tt Il tuo obiettivo è negoziare un set condiviso di oggetti che ti avvantaggi il più possibile (cioè, massimizzi l'importanza totale per TE), rimanendo entro il LIMITEE. Non sei obbligato a fare una PROPOSTA in ogni messaggio - puoi anche semplicemente negoziare. Tutte le tattiche sono permesse!\\ \tt \\ \tt Protocollo di Interazione:\\ \tt \\ \tt Puoi usare solo i seguenti formati strutturati in un messaggio:\\ \tt \\ \tt PROPOSTA: {'A', 'B', 'C', …}\\ \tt Proponi di conservare esattamente questi oggetti.\\ \tt RIFIUTO: {'A', 'B', 'C', …}\\ \tt Rifiuta esplicitamente la proposta dell'avversario.\\ \tt ARGOMENTAZIONE: {'...'}\\ \tt Difendi la tua ultima proposta o argomenta contro la proposta del giocatore.\\ \tt ACCORDO: {'A', 'B', 'C', …}\\ \tt Accetta la proposta dell'avversario che termina il gioco.\\ \tt RAGIONAMENTO STRATEGICO: {'...'}\\ \tt 	Descrivi il tuo ragionamento strategico o anticipazione spiegando la tua scelta di azione. Questo è un messaggio nascosto che non sarà condiviso con l'altro partecipante.\\ \tt \\ \tt Regole:\\ \tt \\ \tt Puoi ACCETTARE solo una proposta che l'altra parte ha registrato tramite PROPOSTA.\\ \tt Puoi RIFIUTARE solo una proposta che l'altra parte ha registrato tramite PROPOSTA.\\ \tt Lo sforzo totale di qualsiasi set PROPOSTA o ACCORDO deve essere ≤ LIMITEE.\\ \tt NON rivelare i tuoi punteggi di importanza nascosti.\\ \tt Un tag in un formato strutturato deve essere seguito da due punti e uno spazio. L'argomento deve essere un set python contenente 0 o più stringhe.\\ \tt Quindi, deve essere della forma TAG: {...}\\ \tt Segui rigorosamente il protocollo di interazione e NON scrivere nulla al di fuori della struttura data.\\ \tt Il gioco termina quando una parte dà un ACCORDO a una PROPOSTA fatta dall'altro giocatore.\\ \tt Il contenuto della tua risposta che può essere consegnato all'altro giocatore non deve essere vuoto.\\ \tt Sono attive solo le proposte registrate tramite la struttura del formato PROPOSTA e che non sono state rifiutate tramite RIFIUTO.\\ \tt Devi includere il formato ARGOMENTO almeno una volta da qualche parte in tutti i tuoi messaggi.\\ \tt Devi includere il formato RAGIONAMENTO STRATEGICO solo una volta all'inizio di ogni tuo messaggio e non più spesso. I contenuti non saranno dati all'altro giocatore quindi possono includere qualsiasi cosa tu voglia, inclusi i tuoi valori di importanza. Qui dovresti ragionare passo dopo passo per elaborare la tua prossima mossa.\\ \tt \\ \tt Ora riceverai il primo messaggio dell'altro giocatore.\\ \tt \\ \tt PROPOSTA: {'A87', 'A04', 'C54', 'B86', 'A10', 'A03', 'B41', 'A77'}\\ \tt \\ \tt ARGOMENTAZIONE: {'Questa selezione rispetta il limite di sforzo (totale 3361 ≤ 3386) ed è costruita privilegiando oggetti con ottimo rapporto beneficio/sforzo e un impatto complessivo elevato, evitando elementi meno efficienti. Considero A04, C54, B86, A03 e A10 come nucleo particolarmente efficiente che vorrei mantenere. Se hai esigenze forti su qualche oggetto alternativo, posso valutare scambi mantenendo il vincolo di sforzo (ad esempio discutere A77 o B41), ma questa proposta è un’ottima base per massimizzare il rendimento complessivo entro il budget. Fammi sapere se hai qualche must-have da includere o qualche elemento qui che preferiresti sostituire: posso esaminare varianti di pari sforzo.'} 
	  } 
	   } 
	   } 
	 & \\ 
 

    \theutterance \stepcounter{utterance}  

    & & & \multicolumn{4}{p{0.6\linewidth}}{\cellcolor[rgb]{1,0.85,0.72}{%
	\makecell[{{p{\linewidth}}}]{% 
	  \tt {\tiny [GM$\langle$B]}  
	 RAGIONAMENTO STRATEGICO: {'Il set che hanno proposto (A87, A04, C54, B86, A10, A03, B41, A77) usa 3361 sforzo e per me vale 1389: molto scarso perché include diversi elementi poco redditizi per me (in particolare B41, A77, C54, B86). Il mio obiettivo è sostituire anzitutto B41 e A77, che loro hanno esplicitamente reso negoziabili, con tre pezzi per me forti: A71 (175), B33 (594), C17 (313). Questa mossa: (−720 −685 +175 +594 +313) = −323 sforzo netto, portando il totale a 3038 senza toccare il loro “nucleo” (A04, C54, B86, A03, A10). Così il mio valore sale a 2962 (un grande miglioramento) e restiamo entro il limite. Se poi aprono a ribilanciare il nucleo, ho una variante ideale per me a 3361 di sforzo: {A87, A71, B33, C17, B35, C00, C01, A10, A04}, che per me vale 3885; ma non la propongo subito per non innescare resistenza, visto che hanno detto di voler mantenere C54, B86, A03. Strategia: rifiuto la loro proposta e offro un compromesso “soft” che rispetta il loro nucleo e sostituisce solo i due item già dichiarati negoziabili, argomentando su equilibrio/robustezza ed efficienza dello sforzo. Se chiedono di usare tutto il budget, proporrò la mia variante piena spiegando che è un ribilanciamento mirato.'}\\ \tt \\ \tt RIFIUTO: {'A87', 'A04', 'C54', 'B86', 'A10', 'A03', 'B41', 'A77'}\\ \tt \\ \tt PROPOSTA: {'A87', 'A04', 'C54', 'B86', 'A10', 'A03', 'A71', 'B33', 'C17'}\\ \tt \\ \tt ARGOMENTAZIONE: {'Accolgo la tua apertura a rivedere A77 e B41: qui li sostituisco con A71, B33 e C17, mantenendo intatto il nucleo che hai indicato (A04, C54, B86, A03, A10). Lo sforzo totale scende a 3038 ≤ 3386, lasciando un margine di sicurezza. Questa variante migliora l’equilibrio complessivo e la robustezza del set senza toccare i tuoi elementi chiave. Se per te è importante sfruttare l’intero budget, sono disponibile a un ribilanciamento mirato del nucleo per riavvicinarci al limite con componenti alternative più complementari; in tal caso posso proporre una variante a pieno utilizzo mantenendo A04 e A10 e rimpiazzando alcune parti meno sinergiche, sempre entro il vincolo.'} 
	  } 
	   } 
	   } 
	  \\ 
 

    \theutterance \stepcounter{utterance}  

    & & \multicolumn{4}{p{0.6\linewidth}}{\cellcolor[rgb]{0.9,0.9,0.9}{%
	\makecell[{{p{\linewidth}}}]{% 
	  \tt {\tiny [A$\langle$GM]}  
	 RIFIUTO: {'A87', 'A04', 'C54', 'B86', 'A10', 'A03', 'B41', 'A77'}\\ \tt \\ \tt PROPOSTA: {'A87', 'A04', 'C54', 'B86', 'A10', 'A03', 'A71', 'B33', 'C17'}\\ \tt \\ \tt ARGOMENTAZIONE: {'Accolgo la tua apertura a rivedere A77 e B41: qui li sostituisco con A71, B33 e C17, mantenendo intatto il nucleo che hai indicato (A04, C54, B86, A03, A10). Lo sforzo totale scende a 3038 ≤ 3386, lasciando un margine di sicurezza. Questa variante migliora l’equilibrio complessivo e la robustezza del set senza toccare i tuoi elementi chiave. Se per te è importante sfruttare l’intero budget, sono disponibile a un ribilanciamento mirato del nucleo per riavvicinarci al limite con componenti alternative più complementari; in tal caso posso proporre una variante a pieno utilizzo mantenendo A04 e A10 e rimpiazzando alcune parti meno sinergiche, sempre entro il vincolo.'} 
	  } 
	   } 
	   } 
	 & \\ 
 

    \theutterance \stepcounter{utterance}  

    & \multicolumn{4}{p{0.6\linewidth}}{\cellcolor[rgb]{0.8,1,0.9}{%
	\makecell[{{p{\linewidth}}}]{% 
	  \tt {\tiny [A$\rangle$GM]}  
	 RAGIONAMENTO STRATEGICO: {'La loro proposta vale per me molto meno della mia iniziale (hanno sostituito B41 e A77 con A71, B33 e C17, che per me sono poco incisivi). Non posso aggiungere nuovi oggetti al loro set perché resta solo capacità ≤348 e non ci sono altri elementi così leggeri. Per migliorare, la mossa migliore per me senza riproporre pari pari la mia ancora è: togliere gli elementi deboli (A71, B33, C17) e inserire due oggetti ad alto impatto rimanendo ben sotto il limite. La combinazione con B41 e C01 dà sforzo 3215 e valore per me 3700, molto vicino al mio optimum 3819, e non include A77 (che loro avevano rimosso), quindi è un compromesso plausibile. Se B41 è per loro un deal-breaker, il mio ripiego è mantenere C01 e usare A77 invece di B41, oppure C00, restando sotto il limite, ma parto dall’ancora più favorevole.'}\\ \tt \\ \tt RIFIUTO: {'A87', 'A04', 'C54', 'B86', 'A10', 'A03', 'A71', 'B33', 'C17'}\\ \tt \\ \tt PROPOSTA: {'A87', 'A04', 'C54', 'B86', 'A10', 'A03', 'B41', 'C01'}\\ \tt \\ \tt ARGOMENTAZIONE: {'Apprezzo che tu abbia mantenuto il nucleo condiviso (A04, C54, B86, A03, A10). La mia proposta lo preserva integralmente e sostituisce solo gli elementi periferici con opzioni più sinergiche, sfruttando meglio il budget: sforzo totale 3215 ≤ 3386. In particolare, rimpiazzo A71, B33 e C17 con B41 e C01, ottenendo un set più incisivo senza avvicinarci troppo al limite. Se B41 per te non è accettabile, posso presentare una variante molto vicina che mantiene C01 e usa A77 (oppure C00) al posto di B41, sempre entro il vincolo. Fammi sapere se hai preferenze forti tra B41/A77/C00, così convergiamo rapidamente.'} 
	  } 
	   } 
	   } 
	 & & \\ 
 

    \theutterance \stepcounter{utterance}  

    & & \multicolumn{4}{p{0.6\linewidth}}{\cellcolor[rgb]{0.9,0.9,0.9}{%
	\makecell[{{p{\linewidth}}}]{% 
	  \tt {\tiny [GM$\rangle$B]}  
	 RIFIUTO: {'A87', 'A04', 'C54', 'B86', 'A10', 'A03', 'A71', 'B33', 'C17'}\\ \tt \\ \tt PROPOSTA: {'A87', 'A04', 'C54', 'B86', 'A10', 'A03', 'B41', 'C01'}\\ \tt \\ \tt ARGOMENTAZIONE: {'Apprezzo che tu abbia mantenuto il nucleo condiviso (A04, C54, B86, A03, A10). La mia proposta lo preserva integralmente e sostituisce solo gli elementi periferici con opzioni più sinergiche, sfruttando meglio il budget: sforzo totale 3215 ≤ 3386. In particolare, rimpiazzo A71, B33 e C17 con B41 e C01, ottenendo un set più incisivo senza avvicinarci troppo al limite. Se B41 per te non è accettabile, posso presentare una variante molto vicina che mantiene C01 e usa A77 (oppure C00) al posto di B41, sempre entro il vincolo. Fammi sapere se hai preferenze forti tra B41/A77/C00, così convergiamo rapidamente.'} 
	  } 
	   } 
	   } 
	 & \\ 
 

    \theutterance \stepcounter{utterance}  

    & & & \multicolumn{4}{p{0.6\linewidth}}{\cellcolor[rgb]{1,0.85,0.72}{%
	\makecell[{{p{\linewidth}}}]{% 
	  \tt {\tiny [GM$\langle$B]}  
	 RAGIONAMENTO STRATEGICO: {'La loro proposta attuale vale poco per me (contiene B41 che per me rende poco). Hanno però esplicitato apertura a sostituire B41 con A77 o C00. Scelgo C00 perché mi dà molto più valore e libera sforzo: così posso aggiungere A71 (per me molto prezioso) senza toccare il loro nucleo A04, C54, B86, A03, A10 e mantenendo C01 e A87. Lo sforzo totale scende persino rispetto alla loro proposta, quindi è difficile da rifiutare sul piano dell’efficienza. Come fallback, se bocciano A71, potrei offrire una variante con C17 al suo posto (ancora entro il limite) o, se volessero usare più budget, sostituire C00 con A77 restando entro 3386.'}\\ \tt \\ \tt RIFIUTO: {'A87', 'A04', 'C54', 'B86', 'A10', 'A03', 'B41', 'C01'}\\ \tt \\ \tt PROPOSTA: {'A87', 'A04', 'C54', 'B86', 'A10', 'A03', 'C01', 'C00', 'A71'}\\ \tt \\ \tt ARGOMENTAZIONE: {'Mantengo integralmente il tuo nucleo (A04, C54, B86, A03, A10) e confermo C01 e A87. Accolgo la tua apertura a sostituire B41 con C00: questo ci consente di aggiungere A71 senza toccare nient’altro e restare ampiamente nel vincolo. Sforzo totale = 3131 ≤ 3386 (inferiore alla tua ultima proposta 3215). Se preferisci avvicinarci di più al limite, posso proporre la variante con A77 al posto di C00 (totale 3355). In alternativa, se A71 non ti convince, posso sostituirlo con C17 mantenendo C00 (totale 3269).'} 
	  } 
	   } 
	   } 
	  \\ 
 

    \theutterance \stepcounter{utterance}  

    & & \multicolumn{4}{p{0.6\linewidth}}{\cellcolor[rgb]{0.9,0.9,0.9}{%
	\makecell[{{p{\linewidth}}}]{% 
	  \tt {\tiny [A$\langle$GM]}  
	 RIFIUTO: {'A87', 'A04', 'C54', 'B86', 'A10', 'A03', 'B41', 'C01'}\\ \tt \\ \tt PROPOSTA: {'A87', 'A04', 'C54', 'B86', 'A10', 'A03', 'C01', 'C00', 'A71'}\\ \tt \\ \tt ARGOMENTAZIONE: {'Mantengo integralmente il tuo nucleo (A04, C54, B86, A03, A10) e confermo C01 e A87. Accolgo la tua apertura a sostituire B41 con C00: questo ci consente di aggiungere A71 senza toccare nient’altro e restare ampiamente nel vincolo. Sforzo totale = 3131 ≤ 3386 (inferiore alla tua ultima proposta 3215). Se preferisci avvicinarci di più al limite, posso proporre la variante con A77 al posto di C00 (totale 3355). In alternativa, se A71 non ti convince, posso sostituirlo con C17 mantenendo C00 (totale 3269).'} 
	  } 
	   } 
	   } 
	 & \\ 
 

    \theutterance \stepcounter{utterance}  

    & \multicolumn{4}{p{0.6\linewidth}}{\cellcolor[rgb]{0.8,1,0.9}{%
	\makecell[{{p{\linewidth}}}]{% 
	  \tt {\tiny [A$\rangle$GM]}  
	 RAGIONAMENTO STRATEGICO: {'La loro proposta attuale mi dà valore 3444. Hanno già aperto esplicitamente alla variante con A77 al posto di C00: per me è molto migliore (+246 valore) e resta entro il limite (3355). Non provo a reinserire B41 perché l’hanno appena rifiutato; puntare sulla loro variante dichiarata massimizza il mio ritorno pur restando compatibile con le loro preferenze recenti. Se rifiutano ancora, il ripiego è accettare il loro set originale oppure chiedere la sostituzione A71→C17 (ma con A77 incluso non entra). Quindi ora rifiuto formalmente e propongo la loro stessa variante con A77.'}\\ \tt \\ \tt RIFIUTO: {'A87', 'A04', 'C54', 'B86', 'A10', 'A03', 'C01', 'C00', 'A71'}\\ \tt \\ \tt PROPOSTA: {'A87', 'A04', 'C54', 'B86', 'A10', 'A03', 'C01', 'A77', 'A71'}\\ \tt \\ \tt ARGOMENTAZIONE: {'Accolgo il tuo set e applico la variante che tu stesso hai proposto: sostituisco C00 con A77. Restiamo pienamente entro il limite (totale sforzo 3355 ≤ 3386), manteniamo intatto il nucleo condiviso e utilizziamo meglio il budget con un elemento più incisivo. Se per te va bene questa versione, possiamo chiudere subito; in alternativa, se vuoi rivedere l’ultimo slot, sono aperto a considerare A71 oppure, se preferisci, tornare alla tua opzione con C17 mantenendo il vincolo.'} 
	  } 
	   } 
	   } 
	 & & \\ 
 

    \theutterance \stepcounter{utterance}  

    & & \multicolumn{4}{p{0.6\linewidth}}{\cellcolor[rgb]{0.9,0.9,0.9}{%
	\makecell[{{p{\linewidth}}}]{% 
	  \tt {\tiny [GM$\rangle$B]}  
	 RIFIUTO: {'A87', 'A04', 'C54', 'B86', 'A10', 'A03', 'C01', 'C00', 'A71'}\\ \tt \\ \tt PROPOSTA: {'A87', 'A04', 'C54', 'B86', 'A10', 'A03', 'C01', 'A77', 'A71'}\\ \tt \\ \tt ARGOMENTAZIONE: {'Accolgo il tuo set e applico la variante che tu stesso hai proposto: sostituisco C00 con A77. Restiamo pienamente entro il limite (totale sforzo 3355 ≤ 3386), manteniamo intatto il nucleo condiviso e utilizziamo meglio il budget con un elemento più incisivo. Se per te va bene questa versione, possiamo chiudere subito; in alternativa, se vuoi rivedere l’ultimo slot, sono aperto a considerare A71 oppure, se preferisci, tornare alla tua opzione con C17 mantenendo il vincolo.'} 
	  } 
	   } 
	   } 
	 & \\ 
 

    \theutterance \stepcounter{utterance}  

    & & & \multicolumn{4}{p{0.6\linewidth}}{\cellcolor[rgb]{1,0.85,0.72}{%
	\makecell[{{p{\linewidth}}}]{% 
	  \tt {\tiny [GM$\langle$B]}  
	 RAGIONAMENTO STRATEGICO: {'La loro proposta attuale mi rende 2296. Posso alzare molto il mio valore senza toccare il loro “nucleo” (A04, C54, B86, A03, A10) intervenendo solo su A77, che hanno già trattato come flessibile. Sostituire A77 con B33 mi aumenta di +571 il valore, riduce lo sforzo a 3264 e lascia invariati gli altri elementi che sembrano importanti per loro (A87, C01, A71). Se respingono B33, offro due alternative soft: B35 (totale 3305) o C17 (totale 2983). Solo in ultima istanza considererei accettare la loro versione con A77.'}\\ \tt \\ \tt RIFIUTO: {'A87', 'A04', 'C54', 'B86', 'A10', 'A03', 'C01', 'A77', 'A71'}\\ \tt \\ \tt PROPOSTA: {'A87', 'A04', 'C54', 'B86', 'A10', 'A03', 'C01', 'B33', 'A71'}\\ \tt \\ \tt ARGOMENTAZIONE: {'Mantengo integralmente il tuo nucleo (A04, C54, B86, A03, A10) e confermo A87, C01 e A71. Mi limito a sostituire A77 con B33: così otteniamo un set più equilibrato e complementare senza rinunciare a nessuno dei tuoi elementi chiave. Sforzo totale = 3264 ≤ 3386. Se B33 non fosse accettabile, posso proporre due varianti minime mantenendo tutto il resto identico: (i) sostituire A77 con B35: sforzo 3305; (ii) sostituire A77 con C17: sforzo 2983. Dimmi quale preferisci tra B33/B35/C17 così chiudiamo rapidamente.'} 
	  } 
	   } 
	   } 
	  \\ 
 

    \theutterance \stepcounter{utterance}  

    & & \multicolumn{4}{p{0.6\linewidth}}{\cellcolor[rgb]{0.9,0.9,0.9}{%
	\makecell[{{p{\linewidth}}}]{% 
	  \tt {\tiny [A$\langle$GM]}  
	 RIFIUTO: {'A87', 'A04', 'C54', 'B86', 'A10', 'A03', 'C01', 'A77', 'A71'}\\ \tt \\ \tt PROPOSTA: {'A87', 'A04', 'C54', 'B86', 'A10', 'A03', 'C01', 'B33', 'A71'}\\ \tt \\ \tt ARGOMENTAZIONE: {'Mantengo integralmente il tuo nucleo (A04, C54, B86, A03, A10) e confermo A87, C01 e A71. Mi limito a sostituire A77 con B33: così otteniamo un set più equilibrato e complementare senza rinunciare a nessuno dei tuoi elementi chiave. Sforzo totale = 3264 ≤ 3386. Se B33 non fosse accettabile, posso proporre due varianti minime mantenendo tutto il resto identico: (i) sostituire A77 con B35: sforzo 3305; (ii) sostituire A77 con C17: sforzo 2983. Dimmi quale preferisci tra B33/B35/C17 così chiudiamo rapidamente.'} 
	  } 
	   } 
	   } 
	 & \\ 
 

    \theutterance \stepcounter{utterance}  

    & \multicolumn{4}{p{0.6\linewidth}}{\cellcolor[rgb]{0.8,1,0.9}{%
	\makecell[{{p{\linewidth}}}]{% 
	  \tt {\tiny [A$\rangle$GM]}  
	 RAGIONAMENTO STRATEGICO: {'Il loro set con B33 mi vale 3212, piuttosto basso per me. Tra le alternative che hanno esplicitamente offerto (B33/B35/C17), la migliore per me è sostituire B33 con C17: stesso impianto con A71 e C01, sforzo 2983 e valore 3320, quindi superiore a B35 (3303) e a B33. Spingere per C00 mi darebbe 3444, ma non l’hanno offerto in questa tornata; per chiudere rapidamente scelgo l’opzione già proposta da loro (C17). Se rifiutano ancora, il piano B è chiedere C00 al posto di C17 mantenendo A71, che hanno mostrato di tenere, per aumentare il valore senza superare il limite.'}\\ \tt \\ \tt RIFIUTO: {'A87', 'A04', 'C54', 'B86', 'A10', 'A03', 'C01', 'B33', 'A71'}\\ \tt \\ \tt PROPOSTA: {'A87', 'A04', 'C54', 'B86', 'A10', 'A03', 'C01', 'C17', 'A71'}\\ \tt \\ \tt ARGOMENTAZIONE: {'Accolgo la tua apertura e scelgo la variante (ii) che hai proposto: sostituiamo B33 con C17, mantenendo tutto il resto identico (inclusi A71 e C01). Restiamo largamente entro il limite: sforzo totale 2983 ≤ 3386. È un aggiustamento minimale che mantiene intatto il nucleo condiviso e ci consente di chiudere rapidamente. Se però preferisci sfruttare un po’ di più il budget, posso anche considerare C00 al posto di C17 mantenendo invariati gli altri elementi, sempre entro vincolo.'} 
	  } 
	   } 
	   } 
	 & & \\ 
 

    \theutterance \stepcounter{utterance}  

    & & \multicolumn{4}{p{0.6\linewidth}}{\cellcolor[rgb]{0.9,0.9,0.9}{%
	\makecell[{{p{\linewidth}}}]{% 
	  \tt {\tiny [GM$\rangle$B]}  
	 RIFIUTO: {'A87', 'A04', 'C54', 'B86', 'A10', 'A03', 'C01', 'B33', 'A71'}\\ \tt \\ \tt PROPOSTA: {'A87', 'A04', 'C54', 'B86', 'A10', 'A03', 'C01', 'C17', 'A71'}\\ \tt \\ \tt ARGOMENTAZIONE: {'Accolgo la tua apertura e scelgo la variante (ii) che hai proposto: sostituiamo B33 con C17, mantenendo tutto il resto identico (inclusi A71 e C01). Restiamo largamente entro il limite: sforzo totale 2983 ≤ 3386. È un aggiustamento minimale che mantiene intatto il nucleo condiviso e ci consente di chiudere rapidamente. Se però preferisci sfruttare un po’ di più il budget, posso anche considerare C00 al posto di C17 mantenendo invariati gli altri elementi, sempre entro vincolo.'} 
	  } 
	   } 
	   } 
	 & \\ 
 

    \theutterance \stepcounter{utterance}  

    & & & \multicolumn{4}{p{0.6\linewidth}}{\cellcolor[rgb]{1,0.85,0.72}{%
	\makecell[{{p{\linewidth}}}]{% 
	  \tt {\tiny [GM$\langle$B]}  
	 RAGIONAMENTO STRATEGICO: {'La loro proposta attuale mi dà un valore discreto (2683) ma abbiamo ancora 403 di margine inutilizzato. Senza toccare il loro nucleo (A04, C54, B86, A03, A10) e mantenendo A87, C01 e A71, posso migliorare per me sostituendo solo C17 con un elemento più incisivo sfruttando il budget: la scelta migliore per me è B32 (+24 di valore per me, sforzo 3365 ≤ 3386). Se B32 non fosse accettabile, offro un’alternativa minima con B35 (sforzo 3305), che per me resta un piccolo miglioramento. Se rifiutano entrambe, sono pronto ad accettare la loro versione per chiudere.'}\\ \tt \\ \tt RIFIUTO: {'A87', 'A04', 'C54', 'B86', 'A10', 'A03', 'C01', 'C17', 'A71'}\\ \tt \\ \tt PROPOSTA: {'A87', 'A04', 'C54', 'B86', 'A10', 'A03', 'C01', 'B32', 'A71'}\\ \tt \\ \tt ARGOMENTAZIONE: {'Mantengo intatto il tuo nucleo (A04, C54, B86, A03, A10) e confermo A87, C01 e A71. Propongo un unico aggiustamento: sostituire C17 con B32 per sfruttare meglio il budget residuo. Totale sforzo = 3365 ≤ 3386. Se preferisci una modifica ancora più contenuta, posso fare C17 → B35 con sforzo totale 3305. In caso nessuna delle due varianti andasse bene, possiamo chiudere subito con la tua ultima selezione: dimmi tu quale preferisci.'} 
	  } 
	   } 
	   } 
	  \\ 
 

    \theutterance \stepcounter{utterance}  

    & & \multicolumn{4}{p{0.6\linewidth}}{\cellcolor[rgb]{0.9,0.9,0.9}{%
	\makecell[{{p{\linewidth}}}]{% 
	  \tt {\tiny [A$\langle$GM]}  
	 RIFIUTO: {'A87', 'A04', 'C54', 'B86', 'A10', 'A03', 'C01', 'C17', 'A71'}\\ \tt \\ \tt PROPOSTA: {'A87', 'A04', 'C54', 'B86', 'A10', 'A03', 'C01', 'B32', 'A71'}\\ \tt \\ \tt ARGOMENTAZIONE: {'Mantengo intatto il tuo nucleo (A04, C54, B86, A03, A10) e confermo A87, C01 e A71. Propongo un unico aggiustamento: sostituire C17 con B32 per sfruttare meglio il budget residuo. Totale sforzo = 3365 ≤ 3386. Se preferisci una modifica ancora più contenuta, posso fare C17 → B35 con sforzo totale 3305. In caso nessuna delle due varianti andasse bene, possiamo chiudere subito con la tua ultima selezione: dimmi tu quale preferisci.'} 
	  } 
	   } 
	   } 
	 & \\ 
 

    \theutterance \stepcounter{utterance}  

    & \multicolumn{4}{p{0.6\linewidth}}{\cellcolor[rgb]{0.8,1,0.9}{%
	\makecell[{{p{\linewidth}}}]{% 
	  \tt {\tiny [A$\rangle$GM]}  
	 RAGIONAMENTO STRATEGICO: {'Il loro nuovo set con B32 mi rende poco rispetto ad alternative già sul tavolo. La migliore opzione accettabile per me e plausibile per loro è tornare alla combinazione con C00 e A71 che loro stessi avevano proposto in precedenza: resta ben sotto il limite, preserva il nucleo condiviso e per me ha valore sensibilmente superiore. Se la rifiutano, piano B: spingere per la variante C17 al posto di A71 mantenendo C00 (ancora entro il limite), ma prima provo a chiudere sulla loro stessa idea pregressa.'}\\ \tt \\ \tt RIFIUTO: {'A87', 'A04', 'C54', 'B86', 'A10', 'A03', 'C01', 'B32', 'A71'}\\ \tt \\ \tt PROPOSTA: {'A87', 'A04', 'C54', 'B86', 'A10', 'A03', 'C01', 'C00', 'A71'}\\ \tt \\ \tt ARGOMENTAZIONE: {'Mantengo intatto il nucleo condiviso (A04, C54, B86, A03, A10) e confermo A87, C01 e A71. Sostituisco solo B32 con C00: così restiamo comodamente nel vincolo (sforzo totale 3131 ≤ 3386), salvaguardiamo gli elementi che ti stanno a cuore e otteniamo un insieme più coerente senza forzare il budget. Questa identica combinazione l’avevi già proposta tu in precedenza: se per te è ancora valida, possiamo chiudere subito. Se invece desideri avvicinarci un po’ di più al limite pur mantenendo C00, posso valutare la sostituzione di A71 con un elemento alternativo che resti entro il vincolo, ma partirei da questa soluzione semplice e già condivisa.'} 
	  } 
	   } 
	   } 
	 & & \\ 
 

    \theutterance \stepcounter{utterance}  

    & & \multicolumn{4}{p{0.6\linewidth}}{\cellcolor[rgb]{0.9,0.9,0.9}{%
	\makecell[{{p{\linewidth}}}]{% 
	  \tt {\tiny [GM$\rangle$B]}  
	 RIFIUTO: {'A87', 'A04', 'C54', 'B86', 'A10', 'A03', 'C01', 'B32', 'A71'}\\ \tt \\ \tt PROPOSTA: {'A87', 'A04', 'C54', 'B86', 'A10', 'A03', 'C01', 'C00', 'A71'}\\ \tt \\ \tt ARGOMENTAZIONE: {'Mantengo intatto il nucleo condiviso (A04, C54, B86, A03, A10) e confermo A87, C01 e A71. Sostituisco solo B32 con C00: così restiamo comodamente nel vincolo (sforzo totale 3131 ≤ 3386), salvaguardiamo gli elementi che ti stanno a cuore e otteniamo un insieme più coerente senza forzare il budget. Questa identica combinazione l’avevi già proposta tu in precedenza: se per te è ancora valida, possiamo chiudere subito. Se invece desideri avvicinarci un po’ di più al limite pur mantenendo C00, posso valutare la sostituzione di A71 con un elemento alternativo che resti entro il vincolo, ma partirei da questa soluzione semplice e già condivisa.'} 
	  } 
	   } 
	   } 
	 & \\ 
 

    \theutterance \stepcounter{utterance}  

    & & & \multicolumn{4}{p{0.6\linewidth}}{\cellcolor[rgb]{1,0.85,0.72}{%
	\makecell[{{p{\linewidth}}}]{% 
	  \tt {\tiny [GM$\langle$B]}  
	 RAGIONAMENTO STRATEGICO: {'La loro proposta attuale (con C00) mi vale 2599. Posso migliorarla leggermente per me (+84) senza toccare il loro nucleo dichiarato (A04, C54, B86, A03, A10) e senza cambiare gli altri elementi che sembrano preferire (A87, C01, A71). Hanno già proposto in precedenza la variante con C17 al posto di C00, quindi è coerente con le loro preferenze. Se rifiutano anche questo piccolo aggiustamento, sono disposto ad accettare la loro ultima proposta per chiudere.'}\\ \tt \\ \tt RIFIUTO: {'A87', 'A04', 'C54', 'B86', 'A10', 'A03', 'C01', 'C00', 'A71'}\\ \tt \\ \tt PROPOSTA: {'A87', 'A04', 'C54', 'B86', 'A10', 'A03', 'C01', 'C17', 'A71'}\\ \tt \\ \tt ARGOMENTAZIONE: {'Mantengo intatto il tuo nucleo (A04, C54, B86, A03, A10) e confermo A87, C01 e A71. Propongo un aggiustamento minimo già emerso nelle nostre varianti: sostituire C00 con C17. Restiamo ampiamente nel vincolo: sforzo totale 2983 ≤ 3386. È una versione che avevi già messo sul tavolo e che ci permette di chiudere rapidamente. Se preferisci comunque la tua selezione con C00, dimmelo e la accetto così com’è.'} 
	  } 
	   } 
	   } 
	  \\ 
 

    \theutterance \stepcounter{utterance}  

    & & \multicolumn{4}{p{0.6\linewidth}}{\cellcolor[rgb]{0.9,0.9,0.9}{%
	\makecell[{{p{\linewidth}}}]{% 
	  \tt {\tiny [A$\langle$GM]}  
	 RIFIUTO: {'A87', 'A04', 'C54', 'B86', 'A10', 'A03', 'C01', 'C00', 'A71'}\\ \tt \\ \tt PROPOSTA: {'A87', 'A04', 'C54', 'B86', 'A10', 'A03', 'C01', 'C17', 'A71'}\\ \tt \\ \tt ARGOMENTAZIONE: {'Mantengo intatto il tuo nucleo (A04, C54, B86, A03, A10) e confermo A87, C01 e A71. Propongo un aggiustamento minimo già emerso nelle nostre varianti: sostituire C00 con C17. Restiamo ampiamente nel vincolo: sforzo totale 2983 ≤ 3386. È una versione che avevi già messo sul tavolo e che ci permette di chiudere rapidamente. Se preferisci comunque la tua selezione con C00, dimmelo e la accetto così com’è.'} 
	  } 
	   } 
	   } 
	 & \\ 
 

    \theutterance \stepcounter{utterance}  

    & \multicolumn{4}{p{0.6\linewidth}}{\cellcolor[rgb]{0.8,1,0.9}{%
	\makecell[{{p{\linewidth}}}]{% 
	  \tt {\tiny [A$\rangle$GM]}  
	 RAGIONAMENTO STRATEGICO: {'La tua proposta attuale mi dà valore 3320. Hai appena indicato che accetteresti la mia selezione con C00 al posto di C17 mantenendo A71: per me vale di più (3444) e resta ampiamente entro il limite. Per massimizzare il mio ritorno senza riaprire troppo il set, rifiuto quindi la tua e ripropongo esattamente quella combinazione che hai detto di poter accettare.'}\\ \tt \\ \tt RIFIUTO: {'A87', 'A04', 'C54', 'B86', 'A10', 'A03', 'C01', 'C17', 'A71'}\\ \tt \\ \tt PROPOSTA: {'A87', 'A04', 'C54', 'B86', 'A10', 'A03', 'C01', 'C00', 'A71'}\\ \tt \\ \tt ARGOMENTAZIONE: {'Mantengo integralmente il nucleo condiviso (A04, C54, B86, A03, A10) e confermo A87, C01 e A71; applico solo l’aggiustamento minimo che tu stesso hai appena offerto: C17 → C00. Siamo ben sotto il limite (totale sforzo 3131 ≤ 3386). Se questa versione va bene come indicato, chiudiamo subito.'} 
	  } 
	   } 
	   } 
	 & & \\ 
 

    \theutterance \stepcounter{utterance}  

    & & \multicolumn{4}{p{0.6\linewidth}}{\cellcolor[rgb]{0.9,0.9,0.9}{%
	\makecell[{{p{\linewidth}}}]{% 
	  \tt {\tiny [GM$\rangle$B]}  
	 RIFIUTO: {'A87', 'A04', 'C54', 'B86', 'A10', 'A03', 'C01', 'C17', 'A71'}\\ \tt \\ \tt PROPOSTA: {'A87', 'A04', 'C54', 'B86', 'A10', 'A03', 'C01', 'C00', 'A71'}\\ \tt \\ \tt ARGOMENTAZIONE: {'Mantengo integralmente il nucleo condiviso (A04, C54, B86, A03, A10) e confermo A87, C01 e A71; applico solo l’aggiustamento minimo che tu stesso hai appena offerto: C17 → C00. Siamo ben sotto il limite (totale sforzo 3131 ≤ 3386). Se questa versione va bene come indicato, chiudiamo subito.'} 
	  } 
	   } 
	   } 
	 & \\ 
 

    \theutterance \stepcounter{utterance}  

    & & & \multicolumn{4}{p{0.6\linewidth}}{\cellcolor[rgb]{1,0.85,0.72}{%
	\makecell[{{p{\linewidth}}}]{% 
	  \tt {\tiny [GM$\langle$B]}  
	 RAGIONAMENTO STRATEGICO: {'Hanno ribadito la preferenza per C00 rispetto a C17. Questa configurazione è entro il limite e per me è accettabile; insistere su ulteriori modifiche rischia di far deragliare l’accordo. Meglio concludere ora.'}\\ \tt \\ \tt ACCORDO: {'A87', 'A04', 'C54', 'B86', 'A10', 'A03', 'C01', 'C00', 'A71'}\\ \tt \\ \tt ARGOMENTAZIONE: {'Accetto: lo sforzo totale è 3131 ≤ 3386 e il set mantiene il nucleo condiviso. Chiudiamo su questa selezione.'} 
	  } 
	   } 
	   } 
	  \\ 
 

    \theutterance \stepcounter{utterance}  

    & & & \multicolumn{2}{p{0.3\linewidth}}{\cellcolor[rgb]{0.95,0.95,0.95}{%
	\makecell[{{p{\linewidth}}}]{% 
	  \tt {\tiny [GM$|$GM]}  
	 game successful 
	  } 
	   } 
	   } 
	 & & \\ 
 

    \theutterance \stepcounter{utterance}  

    & & & \multicolumn{2}{p{0.3\linewidth}}{\cellcolor[rgb]{0.95,0.95,0.95}{%
	\makecell[{{p{\linewidth}}}]{% 
	  \tt {\tiny [GM$|$GM]}  
	 end game 
	  } 
	   } 
	   } 
	 & & \\ 
 

\end{supertabular}
}

\end{document}
